% ==================================================================
%  WEEK 6
% ==================================================================
\chapter{Cloud Security, Privacy, and Trust Management}

\section{Introduction and Week Overview}
Week 6 establishes the theoretical and operational foundations of security in multi-tenant cloud environments. The curriculum systematically progresses through four integrated modules:
\begin{enumerate}[leftmargin=*]
  \item \textbf{Foundational Security Primitives and the CIA Triad:} Mathematical definitions of Confidentiality, Integrity, and Availability; formal models of security policy versus enforcement mechanisms.
  \item \textbf{Attack Taxonomy and Attack Lifecycles:} The four classical structural attacks (Interruption, Interception, Modification, Fabrication); Passive vs Active attacks; reconnaissance, scanning, and exploitation phases.
  \item \textbf{Identity, Authentication, and Authorization:} Deep distinction between Authentication (AuthN) and Authorization (AuthZ); federated protocols (SAML 2.0 and OAuth 2.0); Role-Based Access Control (RBAC).
  \item \textbf{Virtualization Vulnerabilities and Gartner Cloud Risks:} Cross-VM side-channel attacks, cache timing analysis, hypervisor escape, and Gartner's seven cardinal cloud security risks.
\end{enumerate}

\section{Foundational Security Primitives: The CIA Triad}
Security in distributed systems rests upon three core pillars known as the \textbf{CIA Triad}:
\begin{figure}[H]
\centering
\begin{tikzpicture}[node distance=1.8cm, font=\small]
  \node[dark, minimum width=3.8cm, minimum height=1.1cm] (c) {\textbf{Confidentiality}\\(Keeping data \& resources hidden)};
  \node[mid, right=1.5cm of c, minimum width=3.8cm, minimum height=1.1cm] (i) {\textbf{Integrity}\\(Data integrity \& Origin authentication)};
  \node[gold, below=1.2cm of $(c)!0.5!(i)$, minimum width=4.2cm, minimum height=1.1cm] (a) {\textbf{Availability}\\(Timely access to resources)};

  \draw[arr, <->] (c) -- (i);
  \draw[arr, <->] (i) -- (a);
  \draw[arr, <->] (a) -- (c);
\end{tikzpicture}
\caption{The Classical CIA Triad in Cloud Computing.}
\end{figure}

\begin{itemize}[leftmargin=*]
  \item \textbf{Confidentiality:} Preventing the unauthorized disclosure of sensitive information. Cryptographic controls (such as AES-256 for data at rest and TLS 1.3 for data in transit) ensure confidentiality.
  \item \textbf{Integrity:} Guaranteeing that data has not been altered, modified, corrupted, or deleted by unauthorized parties or malfunctioning software. It encompasses:
    \begin{itemize}
      \item \emph{Data Integrity:} State verification ensuring data matches its original format (enforced via SHA-256 cryptographic hashes and HMACs).
      \item \emph{Origin Integrity (Authenticity):} Guaranteeing the genuine origin of the data or service identity.
    \end{itemize}
  \item \textbf{Availability:} Ensuring authorized users have uninterrupted, timely access to computing resources, data stores, and network endpoints without unexpected denial of service.
\end{itemize}

\begin{warnbox}{Integrity Failure Despite Perfect Encryption}
A common misconception is that encryption automatically guarantees data integrity. Encryption scrambles data to preserve \textbf{Confidentiality}. If an administrator or erroneous application software bug alters database entries (e.g., corrupting financial ledger balances), the data remains confidential (still encrypted) but its \textbf{Integrity} is completely destroyed!
\end{warnbox}

\section{Structural Attack Taxonomy}
Prof. Ghosh classifies security violations into four canonical architectural attack patterns based on the source ($S$), destination ($D$), and interceptor ($I$):

\begin{center}
\begin{tabularx}{\textwidth}{>{\bfseries}l l p{3.2cm} X}
\toprule
Attack Type & Violated Property & Structural Mechanism & Concrete Real-World Example\\
\midrule
Interruption & Availability & Attacker destroys or severs the communication path between $S$ and $D$. & Distributed Denial of Service (DDoS), physical fiber cut, RAM starvation.\\
Interception & Confidentiality & Attacker secretly monitors or taps the link, reading packets without altering them. & Packet sniffing via Wireshark, optical fiber tapping, eavesdropping.\\
Modification & Integrity & Attacker intercepts legitimate packets from $S$, alters payload bits, and forwards them to $D$. & Man-in-the-Middle (MitM) payload alteration, SQL injection data corruption.\\
Fabrication & Authenticity & Attacker injects spurious, counterfeit messages into the system claiming to originate from $S$. & DNS spoofing, phishing portals, fake VPN login gateways, ARP poisoning.\\
\bottomrule
\end{tabularx}
\end{center}

\begin{figure}[H]
\centering
\begin{tikzpicture}[node distance=1.5cm, font=\footnotesize]
  % Interruption
  \node[dark] (s1) {S}; \node[dark, right=2.5cm of s1] (d1) {D};
  \draw[arr] (s1) -- node[midway, fill=white] {$\times$} (d1);
  \node[below=0.1cm of s1] {(a) Interruption};

  % Interception
  \node[dark, right=2cm of d1] (s2) {S}; \node[dark, right=2.5cm of s2] (d2) {D};
  \node[mid, below=0.6cm of $(s2)!0.5!(d2)$] (i2) {I};
  \draw[arr] (s2) -- (d2); \draw[arr, dashed] ($(s2)!0.5!(d2)$) -- (i2);
  \node[below=0.8cm of s2] {(b) Interception};

  % Modification
  \node[dark, below=1.8cm of s1] (s3) {S}; \node[dark, right=2.5cm of s3] (d3) {D};
  \node[mid, below=0.6cm of $(s3)!0.5!(d3)$] (i3) {I};
  \draw[arr] (s3) -- (i3); \draw[arr] (i3) -- (d3);
  \node[below=0.8cm of s3] {(c) Modification};

  % Fabrication
  \node[dark, right=2cm of d3] (s4) {S}; \node[dark, right=2.5cm of s4] (d4) {D};
  \node[mid, below=0.6cm of $(s4)!0.5!(d4)$] (i4) {I};
  \draw[arr] (i4) -- (d4);
  \node[below=0.8cm of s4] {(d) Fabrication};
\end{tikzpicture}
\caption{Four Fundamental Classes of Structural Security Attacks.}
\end{figure}

\subsection{Passive vs Active Attacks}
\begin{itemize}[leftmargin=*]
  \item \textbf{Passive Attacks:} The attacker monitors, inspects, or analyzes network transmissions without injecting, modifying, or delaying data packets.
    \begin{itemize}
      \item Examples: Eavesdropping, network packet sniffing, traffic volume/timing analysis.
      \item Operational Implication: Extremely difficult to detect because the system continues functioning normally. Primary defense is proactive \textbf{prevention} via strong end-to-end encryption.
    \end{itemize}
  \item \textbf{Active Attacks:} The attacker actively alters the data stream, creates false streams, disrupts server operation, or modifies system states.
    \begin{itemize}
      \item Examples: Session hijacking, replay attacks, man-in-the-middle tampering, masquerading, Denial of Service (DoS).
      \item Operational Implication: Can be detected via integrity checks, sequence numbers, and anomaly monitoring; focus is on \textbf{detection and rapid automated recovery}.
    \end{itemize}
\end{itemize}

\subsection{The Multi-Phase Attack Lifecycle}
A sophisticated adversary typically executes attacks in structured operational phases:
\begin{enumerate}[leftmargin=*]
  \item \textbf{Reconnaissance (Information Gathering):} Collecting target intelligence without actively exploiting systems.
    \begin{itemize}
      \item \emph{Passive Reconnaissance:} Scraping public databases, social media, WHOIS records, and DNS records.
      \item \emph{Active Reconnaissance:} DNS zone transfers, ping sweeps, and network service mapping.
    \end{itemize}
  \item \textbf{Scanning and Probing:} Using port scanners (e.g., Nmap) and vulnerability scanners to discover listening network services, operating system versions, and known unpatched CVE vulnerabilities.
  \item \textbf{Gaining Access (Exploitation):} Executing targeted exploit payloads (such as buffer overflow attacks, SQL injections, or remote code execution) to bypass security boundaries and achieve administrative privileges.
  \item \textbf{Maintaining Access and Persistence:} Installing stealth rootkits, trojans, or backdoor SSH keys to retain control across server reboots.
  \item \textbf{Covering Tracks:} Altering or deleting system log files (\texttt{/var/log/auth.log}), resetting file modification timestamps, and disabling intrusion detection alerts.
\end{enumerate}

\section{Identity Management: Authentication vs Authorization}
A central security theme emphasized in Week 6 is the crisp formal distinction between Authentication and Authorization:

\begin{center}
\begin{tabularx}{\textwidth}{>{\bfseries}l p{6.5cm} X}
\toprule
Dimension & Authentication (AuthN) & Authorization (AuthZ)\\
\midrule
Core Question & ``Who are you?'' (Verification of Identity) & ``What actions are you permitted to execute?''\\
Mechanism & Verification of digital credentials (passwords, cryptographic tokens, biometric signatures, MFA). & Evaluation of permissions, Access Control Lists (ACLs), or Role-Based policies (RBAC).\\
Timing & Occurs \textbf{first} at system ingress before session establishment. & Evaluated \textbf{continuously} whenever an authenticated subject attempts an operation.\\
Protocol Standards & Kerberos, OpenID Connect (OIDC), SAML 2.0 Identity Assertion, FIDO2 / WebAuthn. & OAuth 2.0 (access tokens and scopes), XACML, AWS IAM Policy documents.\\
Failure Result & HTTP 401 Unauthorized (Unauthenticated identity). & HTTP 403 Forbidden (Identity verified, but lacks sufficient permissions).\\
\bottomrule
\end{tabularx}
\end{center}

\begin{defbox}{Security Assertion Markup Language (SAML 2.0) vs OAuth 2.0}
\textbf{SAML 2.0:} An XML-based federated protocol designed primarily for cross-domain Web Single Sign-On (\textbf{Authentication}). An Identity Provider (IdP) signs an XML assertion containing identity attributes and transmits it to a Service Provider (SP) via the user's browser.\\
\textbf{OAuth 2.0:} A lightweight, JSON-based framework designed for delegated \textbf{Authorization}. It issues short-lived Access Tokens (such as JSON Web Tokens, JWTs) granting a third-party client scoped access to a user's protected resources on a Resource Server without revealing user credentials.
\end{defbox}

\section{Virtualization Security and Multi-Tenancy Exploits}
\subsection{Hypervisor Isolation and Virtual Machine Escape}
In a public cloud, multiple untrusted virtual machines co-exist on the same physical bare-metal host. The hypervisor (Virtual Machine Monitor, VMM) is responsible for enforcing strict isolation:
\begin{itemize}[leftmargin=*]
  \item \textbf{VM Escape:} The most catastrophic virtualization exploit, where code running inside a guest operating system breaks out of the virtualized sandbox and interacts directly with the host hypervisor or host OS kernel. Once escaped, the attacker acquires total control over all co-located tenant VMs on that physical server.
\end{itemize}

\subsection{Cross-VM Side-Channel Attacks (Ristenpart et al.)}
In their seminal 2009 research paper (\emph{Hey, You, Get Off of My Cloud: Exploring Information Leakage in Third-Party Compute Clouds}), Ristenpart et al. demonstrated that physical resource sharing in public clouds (specifically Amazon EC2) introduces subtle side-channel vulnerabilities:

\begin{enumerate}[leftmargin=*]
  \item \textbf{Co-Residence Verification:} An attacker probes cloud IP ranges to identify the target victim's network address. By exploiting internal network topologies (comparing Dom0 IP hops and round-trip ping times $\le 1\text{ ms}$), the attacker verifies that their own VM resides on the identical physical hardware as the victim.
  \item \textbf{Shared Hardware Contention:} Even when CPU registers and memory pages are strictly isolated by hardware MMU page tables, physical cores share hardware caches:
    \begin{itemize}
      \item L1 / L2 caches (per core) and unified \textbf{L3 Last Level Cache (LLC)}.
      \item Memory buses, CPU branch target buffers (BTB), and hyperthreading pipelines.
    \end{itemize}
  \item \textbf{Cache-Timing Side Channels (Prime + Probe):}
    \begin{itemize}
      \item \emph{Prime Phase:} The attacker fills specific cache sets with their own data lines.
      \item \emph{Idle Window:} The attacker waits while the victim VM executes cryptographic operations (such as modular exponentiation during RSA signing).
      \item \emph{Probe Phase:} The attacker re-reads their own cached memory addresses, timing the access latency. Memory addresses that were evicted by the victim take longer to reload from main RAM.
      \item \emph{Key Extraction:} By observing which cache sets experienced evictions, the attacker reconstructs internal CPU access patterns, leaking private cryptographic keys without ever violating memory permissions!
    \end{itemize}
\end{enumerate}

\section{Gartner's Seven Cardinal Cloud Security Risks}
The lecture highlights the influential risk taxonomy published by Gartner for enterprise cloud computing adopters:
\begin{enumerate}[leftmargin=*]
  \item \textbf{Privileged User Access:} Cloud providers maintain system administrators with root access to physical hosts and virtualization hypervisors. Organizations must demand strict administrative auditing, background vetting, and cryptographic envelope encryption.
  \item \textbf{Regulatory Compliance:} Enterprises remain legally liable for data compliance (HIPAA, PCI-DSS, ISO 27001) even when data is hosted externally. Providers must submit to external third-party compliance audits.
  \item \textbf{Data Location and Sovereignty:} Cloud consumers often cannot specify the exact physical storage medium or geographic jurisdiction where data is stored. Trans-border data movement may violate national data residency laws (e.g., GDPR in the European Union or local data sovereignty mandates).
  \item \textbf{Data Segregation / Isolation:} In shared multi-tenant databases and SAN volumes, data from competing corporations resides on adjacent disk sectors. Multi-tenant encryption schemes must guarantee complete logical segregation.
  \item \textbf{Recovery and Business Continuity:} In the event of catastrophic physical outages or ransomware events at the provider's facility, providers must have verified Recovery Point Objectives (RPO) and Recovery Time Objectives (RTO).
  \item \textbf{Investigative Support (E-Discovery):} Conducting forensic investigations or digital discovery in a multi-tenant cloud is exceptionally difficult without violating the privacy rights of other co-located tenants.
  \item \textbf{Long-Term Viability:} What happens to an organization's mission-critical data if the cloud provider goes bankrupt, gets acquired by a hostile competitor, or abruptly terminates service? Data portability and non-proprietary exit strategies are essential.
\end{enumerate}

\section{Comprehensive Solved Assignment Problems (Week 6)}

\begin{solved}{Database Bug and Security Objective Failure (Assignment 6, Q1)}
\Given A database stores encrypted customer records. An administrator erroneously alters account balances due to an unhandled application bug. Identify the security objective that failed.
\Steps
\begin{enumerate}[leftmargin=*]
  \item Encryption preserves \textbf{Confidentiality} (unauthorized parties cannot read the records).
  \item Availability is intact because users can still query the system.
  \item Erroneous or unauthorized modification of data directly compromises \textbf{Integrity}.
\end{enumerate}
\Answer \textbf{Option B: Integrity only}.
\end{solved}

\begin{solved}{Fake VPN Portal Exploitation (Assignment 6, Q2)}
\Given An employee visits a fake VPN portal that visually imitates the legitimate corporate login page. Before any credentials are typed, which security property has failed?
\Steps
\begin{enumerate}[leftmargin=*]
  \item No confidential password has been disclosed yet.
  \item The fraudulent portal is masquerading as the genuine corporate gateway (a fabrication attack).
  \item The security property that is broken is \textbf{Authenticity} (the claim of genuine identity).
\end{enumerate}
\Answer \textbf{Option D: Authenticity}.
\end{solved}

\begin{solved}{Classification of Active vs Passive Attacks (Assignment 6, Q3)}
\Given Identify which attack is NOT considered a passive attack among: Packet sniffing, Traffic analysis, Session hijacking, Network monitoring.
\Steps
\begin{enumerate}[leftmargin=*]
  \item Packet sniffing, traffic analysis, and network monitoring merely inspect transmission without altering packets (Passive).
  \item \textbf{Session hijacking} actively injects malicious packets, alters sequence numbers, or terminates the user's connection to seize control of the session (Active).
\end{enumerate}
\Answer \textbf{Option C: Session hijacking}.
\end{solved}

\begin{solved}{Inability to Modify Transmitted Payloads (Assignment 6, Q4)}
\Given Which attack cannot directly modify the contents of transmitted data?
\Steps
\begin{enumerate}[leftmargin=*]
  \item Replay attacks inject recorded packets; MitM attacks alter bits in flight; Masquerading sends forged data.
  \item \textbf{Eavesdropping} is purely passive listening and cannot modify data contents.
\end{enumerate}
\Answer \textbf{Option C: Eavesdropping}.
\end{solved}

\begin{solved}{Denial of Service (DoS) Characteristics (Assignment 6, Q5)}
\Given What best describes a Denial-of-Service attack?
\Steps
\begin{enumerate}[leftmargin=*]
  \item An adversary floods servers or networks with overwhelming request volumes, consuming compute/bandwidth resources so legitimate users cannot access the service.
\end{enumerate}
\Answer \textbf{Option B: An attacker floods a server with excessive requests, making the service unavailable to legitimate users}.
\end{solved}

\begin{solved}{Reconnaissance vs Exploitation Phase (Assignment 6, Q6)}
\Given Identify the activity least likely to be considered part of the reconnaissance phase.
\Steps
\begin{enumerate}[leftmargin=*]
  \item DNS enumeration, port scanning, and searching public websites are reconnaissance activities.
  \item \textbf{Exploiting a buffer overflow} to gain root access is part of the active \emph{Exploitation} phase, not reconnaissance.
\end{enumerate}
\Answer \textbf{Option C: Exploiting a buffer overflow vulnerability to gain administrator privileges}.
\end{solved}

\begin{solved}{Authentication vs Authorization Statements (Assignment 6, Q7)}
\Given Evaluate:
\begin{itemize}
  \item Statement I: Authentication determines whether a user is allowed to access a resource.
  \item Statement II: Authorization determines what actions an authenticated user is permitted to perform.
\end{itemize}
\Steps
\begin{enumerate}[leftmargin=*]
  \item Statement I is FALSE: Authentication verifies \emph{identity} (who the user is), not whether they are allowed to access a specific resource.
  \item Statement II is TRUE: Authorization evaluates permissions to determine permitted actions.
\end{enumerate}
\Answer \textbf{Statement I is False, and Statement II is True}.
\end{solved}

\begin{solved}{Data Location and Sovereignty Concerns (Assignment 6, Q8)}
\Given A multinational corporation stores EU citizen records in a US-based cloud data center and faces legal scrutiny under GDPR. Which Gartner cloud risk applies?
\Steps
\begin{enumerate}[leftmargin=*]
  \item European privacy mandates require personal citizen data to reside within compliant jurisdictions.
  \item Moving data across geographic borders directly triggers \textbf{Data Location and Sovereignty} risk.
\end{enumerate}
\Answer \textbf{Option C: Data Location}.
\end{solved}

\begin{solved}{Multi-Tenant Cache-Timing Exploits (Assignment 6, Q9)}
\Given How does an attacker extract private cryptographic keys in a cross-VM side-channel attack?
\Steps
\begin{enumerate}[leftmargin=*]
  \item By establishing co-residence on the same physical CPU socket and monitoring access latencies to shared L3 CPU cache lines (Prime+Probe technique).
\end{enumerate}
\Answer \textbf{Option B: By measuring CPU cache access latency on shared physical hardware}.
\end{solved}

\begin{solved}{Hypervisor Sandbox Escape Threat (Assignment 6, Q10)}
\Given What is the security impact of a successful VM Escape vulnerability?
\Steps
\begin{enumerate}[leftmargin=*]
  \item The attacker breaks through guest OS virtualization barriers to execute arbitrary code with host hypervisor privileges, compromising all co-located tenant VMs.
\end{enumerate}
\Answer \textbf{Option A: The attacker gains direct access to the host hypervisor, compromising all tenant VMs on the physical host}.
\end{solved}
